\chapter{Literature Review}\label{ch:background}

\section{Building Energy Flexibility and CityLearn}

Reinforcement learning for building energy management has moved from isolated
HVAC control toward coordinated control of flexible communities. The motivation
is clear: buildings consume a large share of electricity and can provide demand
flexibility, but controllers must balance energy, comfort, cost, carbon, and
operational constraints. Recent surveys emphasize that RL is promising for
adaptive building control, while also noting practical barriers such as safety,
sample efficiency, sim-to-real transfer, and benchmark comparability
\cite{zandi2023tenquestions,safeDRL2025}.

CityLearn directly addresses the benchmark problem. CityLearn v1 introduced an
OpenAI Gym-style environment for demand response using deep RL
\cite{citylearn2019}. The later CityLearn work standardized multi-agent RL
research for urban demand response and made comparison between algorithms more
reproducible \cite{citylearn2020}. CityLearn v2 extends the environment with
energy-flexible, resilient, occupant-centric, and carbon-aware community
control, including richer distributed energy resource models and updated
evaluation use cases \cite{citylearnv2}. This makes CityLearn an appropriate
environment for this project because the research question is not simply
whether an agent can reduce energy cost, but whether multiple building agents
can coordinate under shared portfolio-level objectives.

\section{Annex 96 CE1 Experimental Setting}

This project uses Annex 96 Common Exercise 1 (CE1) as a benchmark setting in
CityLearn. Each CE1 schema defines 25 single-family homes, one weather file, one
district target file, and one energy-simulation file per home. The district
target file provides the reference load that the whole portfolio should follow,
while the per-building files provide local demand, indoor temperature, solar
generation, and other time-series observations. The important point for this
project is that CE1 is a portfolio load-tracking task, not a physical network
control task.

The CE1 schemas used here do not provide a feeder map, floor-plan distribution,
geographic coordinates, or an adjacency list between homes. Therefore the
current grouping method cannot use physical proximity or electrical topology.
Instead, buildings are grouped using flexibility-related device features that
are available after CityLearn initializes the environment, mainly electrical
storage capacity and total HVAC nominal power. This matches the K-means
implementation used in the experiments.

Another data limitation is pricing and carbon intensity. The CE1 schema contains
observation names for electricity pricing and carbon intensity, but the VT and TX
CE1 building entries used in this project set both the pricing file and the
carbon-intensity file to null. CityLearn then fills these series with zeros.
As a result, cost and carbon-emission metrics still have calculation formulas in
CityLearn, but they are zero or null in the current experiments and should not be
interpreted as meaningful performance outcomes.

\section{MARL for Energy Systems}

Multi-agent reinforcement learning is a natural formulation for energy
communities because each building, inverter, or household can be treated as an
agent. General MARL surveys highlight three recurring difficulties:
non-stationarity from concurrently learning agents, partial observability, and
credit assignment under shared rewards \cite{marlOverview2019}. In energy
systems these issues are amplified by heterogeneous devices, different comfort
states, network constraints, and long temporal dependencies.

Recent energy-system studies show why MARL is appropriate but also why
structure is necessary. PowerGridworld provides a Gym-style environment for
multi-agent power-system control and demonstrates MARL with PPO and MADDPG
\cite{powergridworld2022}. Charbonnier et al. study scalable residential
energy flexibility and show that centralized-but-factored critics can improve
coordination and scalability for decentralized homes \cite{charbonnier2025}.
Physics-informed MARL for distribution-network voltage control further shows
that graph and physical structure can improve robustness and sample efficiency
\cite{physicsMADRL2024}. These works support the design principle used in this
project: scalable energy MARL should not treat agents as an unstructured list;
it should use building similarity, graph structure, or communication structure.

\section{Rule-Based and Single-Agent Baselines}

Rule-based control (RBC) is an important non-learning baseline in CityLearn.
An RBC controller uses fixed rules, often based on hour-of-day schedules, to
charge or discharge storage and to set heating or cooling device availability.
It is simple, transparent, and cheap to run, but it cannot adapt its policy to
the CE1 reference profile or to interactions between buildings. In this project,
RBC is therefore used as a practical reference point: a MARL method should
eventually do more than repeat a fixed schedule.

Independent single-agent RL baselines are also important. PPO is a robust
on-policy policy-gradient method that limits destructive policy updates through
a clipped surrogate objective \cite{ppo2017}. It is often used as a stable
baseline because it is comparatively simple and reliable. SAC is an off-policy
maximum-entropy actor-critic method that learns stochastic policies and
encourages exploration through entropy regularization \cite{sac2018,
sacApplications2018}. In this project, independent PPO and SAC provide
baseline controllers that answer a basic question: how much performance can be
obtained without explicit inter-building coordination?

On-policy and off-policy methods also differ in how they use data. On-policy
methods such as PPO and MAPPO collect trajectories using the current policy and
then update from those fresh trajectories. This is usually stable and easier to
control, but it uses data less efficiently. Off-policy methods such as SAC can
store past experience in a replay buffer and reuse it for many updates. This can
be more sample-efficient and exploratory, but it can also be more sensitive to
distribution shift between old data and the current policy.

However, independent baselines are expected to struggle with portfolio-level
tracking. Each building receives only local observations and cannot know how
other buildings will respond at the same time step. This motivates moving from
independent learning to CTDE and communication-aware MARL.

\section{Centralized Training with Decentralized Execution}

CTDE is the dominant framework for cooperative MARL. MADDPG uses centralized
critics that condition on other agents' actions to reduce non-stationarity in
actor-critic learning \cite{maddpg2017}. COMA addresses credit assignment by
using a centralized critic with a counterfactual baseline that estimates each
agent's marginal contribution \cite{coma2018}. Value-decomposition approaches
such as VDN and QMIX decompose a team value function into per-agent components.
QMIX adds a monotonic mixing constraint so that decentralized greedy actions
remain consistent with the centralized value function \cite{vdn2017,qmix2018}.

MAPPO is a policy-gradient CTDE method that adapts PPO to cooperative
multi-agent settings. Yu et al. show that, with careful implementation,
PPO-style on-policy learning is surprisingly effective in cooperative
multi-agent games \cite{mappo2022}. MAPPO is a good match for CityLearn because
CityLearn has a continuous bounded action space, and the portfolio objective is
cooperative. This project therefore uses MAPPO as the main learning backbone and
reuses the existing on-policy implementation rather than reimplementing the
optimizer.

\section{Parameter Sharing, Grouping, and K-Means}

Parameter sharing is a common technique for scaling MARL: multiple agents use
the same actor network, reducing the number of trainable parameters and
improving sample efficiency. The limitation is that full parameter sharing can
be harmful when agents are heterogeneous. Selective parameter sharing addresses
this by partitioning agents according to abilities or goals, then sharing
parameters within each partition \cite{seps2021}. This idea is directly
relevant to the CityLearn portfolio, where buildings differ in battery capacity,
thermal flexibility, and load profile.

The grouping method used in this project is K-means clustering. K-means
partitions data into $K$ sets by minimizing within-cluster squared distances to
cluster centroids \cite{kmeans1967}. In this implementation, buildings are
clustered using flexibility-related features such as battery storage capacity
and aggregate HVAC capacity. The result is a grouped actor architecture: agents
within the same cluster share an actor, while the full portfolio shares a
centralized critic. This provides a concrete link between classical clustering,
selective parameter sharing, and building-energy heterogeneity.

\section{Communication in MARL}

Communication is another way to overcome partial observability and
non-stationarity. A recent survey of multi-agent deep RL with communication
argues that communication broadens agents' effective observations and supports
cooperation, but also introduces design choices about message content, receiver
selection, communication frequency, and overhead \cite{commSurvey2024}. This is
especially relevant for building portfolios: all-to-all communication may be
unnecessary or noisy, while no communication may be too weak for aggregate
tracking.

Several communication methods guide this project. DIAL trains differentiable
communication channels under centralized learning and decentralized execution
\cite{dial2016}. CommNet performs differentiable message passing by averaging
hidden states from other agents and learning a neural update
\cite{commnet2016}. TarMAC replaces uniform broadcasting with attention-based
targeted communication so that agents can learn whom to attend to
\cite{tarmac2019}. PowerNet, developed for scalable power-grid control, uses
local-neighbor communication and is closely aligned with the idea that energy
systems have useful graph structure \cite{powernet2020}.

The communication modules in this project combine these ideas with the
CityLearn portfolio structure. CommNet tests simple mean-pooling communication.
PowerNet-style modules test more structured information passing: each agent
keeps its own encoded feature, aggregates selected peer information, concatenates
the two streams, and applies a residual update so the original local information
is not overwritten. TarMAC tests learned attention. The weighted
same-group/other-group module is a project-specific middle ground: it gives
higher weight to messages from similar buildings and lower weight to global
context. This method is intended to preserve scalability while still allowing
the policy to coordinate at portfolio level.

\section{Research Gap}

The literature shows that RL for building control is active, CityLearn provides
a strong benchmark, CTDE methods are effective for cooperative MARL, and
communication can improve coordination. The gap addressed here is the
combination of these ideas for CityLearn portfolio load tracking:

\begin{itemize}
  \item most CityLearn controllers do not explicitly compare multiple
  structured communication protocols;

This project therefore focuses on grouped MAPPO with structured communication,
using CityLearn and Annex 96 metrics as the application-grounded evaluation
setting.
